\documentclass[../../main.tex]{subfiles}
\begin{document}
\section{叶状结构}\label{section:ch19-foliations}

当我们把光滑流形 $M$ 上一个对合的秩-$k$ 分布的所有极大积分流形汇集到一起时，就得到 $M$ 的一个分解，它由若干 $k$ 维子流形组成，这些子流形在局部像平坦坐标图中的切片那样"拼合"在一起。

为了更精确地表达"拼合在一起"的含义，需要把平坦坐标图的概念稍加推广。设 $M$ 是光滑 $n$ 维流形，$\mathscr F$ 是 $M$ 的任意一族 $k$ 维子流形。若 $M$ 的一个光滑坐标图 $(U,\varphi)$ 满足 $\varphi(U)$ 是 $\mathbb{R}^n$ 中的一个立方体，且 $\mathscr F$ 中每个子流形与 $U$ 的交集要么为空，要么是形如 $x^{k+1}=c^{k+1},\dots,x^n=c^n$ 的 $k$ 维切片的可数并，则称 $(U,\varphi)$ 对 $\mathscr F$ 是\emph{平坦的}。我们把 $M$ 上一个\emph{$k$ 维叶状结构}定义为 $M$ 的一族互不相交、连通、非空的浸入 $k$ 维子流形（称为该叶状结构的\emph{叶}），其并集为 $M$，并且在每一点 $p\in M$ 的一个邻域内都存在对 $\mathscr F$ 平坦的坐标图。

\begin{example}[叶状结构]\label{example:ch19-foliations}
(a) $\mathbb{R}^n$ 中所有平行于 $\mathbb{R}^k\times\{0\}$ 的 $k$ 维仿射子空间构成 $\mathbb{R}^n$ 的一个 $k$ 维叶状结构。
(b) 当 $x$ 遍历 $\mathbb{R}^n\setminus\{0\}$ 时，形如 $\{tx\colon t>0\}$ 的开射线族构成 $\mathbb{R}^n\setminus\{0\}$ 的一个 $1$ 维叶状结构。
(c) 以 $0$ 为中心的所有球面构成 $\mathbb{R}^n\setminus\{0\}$ 的一个 $(n-1)$ 维叶状结构。
(d) 若 $M$ 与 $N$ 是连通光滑流形，则当 $q$ 遍历 $N$ 中的点时，形如 $M\times\{q\}$ 的子集族构成 $M\times N$ 的一个叶状结构，其每一叶都与 $M$ 微分同胚。例如，当 $q\in S^1$ 时，形如 $S^1\times\{q\}\subseteq T^2$ 的所有圆周族给出环面 $T^2$ 的一个叶状结构（见图\cref{figure:ch19-torus-foliation}(a)）。$T^2$ 的另一个叶状结构由形如 $\{p\}\times S^1$ 的圆周族给出（见图\cref{figure:ch19-torus-foliation}(b)）。
(e) 若 $\alpha$ 是一个固定实数，则当 $\theta$ 遍历 $\mathbb{R}$ 时，形如
\begin{align*}
\gamma(t)=\bigl(e^{\mathrm{i}t},e^{\mathrm{i}(\alpha t+\theta)}\bigr)
\end{align*}
的所有曲线的像构成环面的一个 $1$ 维叶状结构（见图\cref{figure:ch19-torus-foliation}(c)）。若 $\alpha$ 是有理数，则每一叶都是嵌入圆周；而若 $\alpha$ 是无理数，则每一叶都是稠密的（见例\cref{example:ch04-20} 与习题 4-4）。
(f) 由下列方程在 $(y,z)$ 平面上定义的曲线的连通分支构成 $\mathbb{R}^2$ 的一个叶状结构（见图\cref{figure:ch19-r2-r3-foliation}(a)）：
\begin{align*}
z=\sec y+c,\quad c\in\mathbb{R};\qquad y\ne k+\tfrac12,\quad k\in\mathbb{Z}.
\end{align*}
(g) 若把上一例中的曲线绕 $z$ 轴旋转，则得到 $\mathbb{R}^3$ 的一个 $2$ 维叶状结构，其中有些叶与圆盘微分同胚，有些与圆柱微分同胚（见图\cref{figure:ch19-r2-r3-foliation}(b)）。
\end{example}

\begin{figure}[ht]
\centering
\includegraphics[width=0.8\textwidth]{fig-19-5.png}
\caption{环面的叶状结构}\label{figure:ch19-torus-foliation}
\end{figure}

\begin{figure}[ht]
\centering
\includegraphics[width=0.8\textwidth]{fig-19-6.png}
\caption{$\mathbb{R}^2$ 和 $\mathbb{R}^3$ 的叶状结构}\label{figure:ch19-r2-r3-foliation}
\end{figure}

关于叶状结构的主要事实是：叶状结构与对合分布一一对应。其中一个方向由下一个命题给出，它是定义的一个简单推论。

\begin{proposition}\label{proposition:ch19-叶给出对合分布}
设 $\mathscr F$ 是光滑流形 $M$ 上的一个叶状结构。$\mathscr F$ 的各叶的切空间构成的集合是 $M$ 上的一个对合分布。
\end{proposition}

\begin{exercise}\label{exercise:ch19-20}
证明命题\cref{proposition:ch19-叶给出对合分布}。
\end{exercise}

Frobenius 定理使我们能够得出下列逆命题，它要深刻得多。顺便指出，值得注意的是这个结果是引入浸入子流形概念的两个主要原因之一（另一个是为了研究 Lie 子群）。

\begin{theorem}[整体 Frobenius 定理]\label{theorem:ch19-整体frobenius}
设 $D$ 是光滑流形 $M$ 上的一个对合分布。$D$ 的所有极大连通积分流形构成的族是 $M$ 的一个叶状结构。
\end{theorem}

该定理是下列引理的一个简单推论。

\begin{lemma}\label{lemma:ch19-积分流形并}
设 $D\subseteq TM$ 是一个对合分布，$\{N_\alpha\}_{\alpha\in A}$ 是 $D$ 的任意一族有公共点的连通积分流形。则 $N=\bigcup_\alpha N_\alpha$ 具有唯一的光滑流形结构，使它成为 $D$ 的一个连通积分流形。
\end{lemma}

\begin{proof}
若能构造出使 $N$ 成为 $D$ 的积分流形的一个拓扑与光滑流形结构，则定理\cref{theorem:ch05-33} 表明该拓扑与光滑结构是唯一确定的，因为积分流形是弱嵌入的。

为构造拓扑，首先需要证明对每个 $\alpha,\beta\in A$，$N_\alpha\cap N_\beta$ 在 $N_\alpha$ 与 $N_\beta$ 中都是开的。任取 $q\in N_\alpha\cap N_\beta$，在 $q$ 的一个邻域 $W$ 上选取 $D$ 的一个平坦坐标图（见图\cref{figure:ch19-union}）。设 $V_\alpha,V_\beta$ 分别表示 $N_\alpha\cap W$、$N_\beta\cap W$ 中含 $q$ 的分支。由命题\cref{proposition:ch19-局部结构}，$V_\alpha$ 与 $V_\beta$ 都是单个切片在子空间拓扑下的开子集；由于二者都含 $q$，它们必都含于同一切片 $S$ 中。于是 $V_\alpha\cap V_\beta$ 在 $S$ 中是开的，在 $N_\alpha$ 与 $N_\beta$ 中也都是开的，故 $q$ 在 $N_\alpha$ 与 $N_\beta$ 中各有一个含于 $N_\alpha\cap N_\beta$ 的邻域。

规定 $N$ 的一个子集 $U\subseteq N$ 为开当且仅当对每个 $\alpha$，$U\cap N_\alpha$ 在 $N_\alpha$ 中是开的，由此在 $N$ 上定义一个拓扑。利用上一段的结果，容易验证这是一个拓扑，且每个 $N_\alpha$ 都是 $N$ 的一个开子空间。在该拓扑下，$N$ 是 $k$ 维局部 Euclid 的，因为每个 $q\in N$ 都在某个 $N_\alpha$ 中有一个坐标邻域 $V$，而由于 $N_\alpha$ 在 $N$ 中开，$V$ 在 $N$ 中也是开的。此外，包含映射 $N\hookrightarrow M$ 连续：对任意开子集 $U\subseteq M$，由于对每个 $\alpha$，$U\cap N_\alpha$ 在 $N_\alpha$ 中开，故 $U\cap N$ 在 $N$ 中开。

为看出 $N$ 是 Hausdorff 的，令 $q,q'$ 为 $N$ 中不同两点。存在 $M$ 中分别含 $q$、$q'$ 的互不相交开子集 $U,U'$；由于包含 $N\hookrightarrow M$ 连续，$N\cap U$ 与 $N\cap U'$ 是 $N$ 中分别含 $q$、$q'$ 的互不相交开子集。

接着证明 $N$ 是第二可数的。可用可数多个 $D$ 的平坦坐标图 $\{W_i\}$ 覆盖 $M$。只需证明对每个 $i$，$N\cap W_i$ 含于可数多个切片的并中即可，因为单个切片的任意开子集都是第二可数的，从而 $N$ 可表示为可数多个子集的并，其中每一个在 $N$ 中都是第二可数且开的。

令 $p_0$ 为含于每个 $N_\alpha$ 中的一点。称某个 $W_k$ 的一个切片 $S$ 是从 $p_0$ \emph{可达的}，若存在有限的指标序列 $i_1,\dots,i_m$，并且对每个 $i_j$ 有切片 $S_{i_j}\subseteq W_{i_j}$，满足 $p_0\in S_{i_1}$、$S_{i_m}=S$，以及对每个 $j=1,\dots,m-1$，$S_{i_j}\cap S_{i_{j+1}}\ne\varnothing$（见图\cref{figure:ch19-accessible}）。

设 $W_k$ 是可数平坦坐标图之一，并设 $S\subseteq W_k$ 是含某点 $q\in N$ 的一个切片。则 $q$ 含于某个积分流形 $N_\alpha$ 中。由于 $p_0$ 也在 $N_\alpha$ 中，存在连续路径 $\gamma\colon[0,1]\to N_\alpha$ 连接 $p_0$ 与 $q$。由于 $\gamma([0,1])$ 紧，存在有限多个数 $0=t_0<t_1<\cdots<t_m=1$，使得对每个 $j=1,\dots,m$，集合 $\gamma([t_{j-1},t_j])$ 含于某个平坦坐标图 $W_{i_j}$ 中。由于 $\gamma([t_{j-1},t_j])$ 连通，它含于 $W_{i_j}\cap N_\alpha$ 的单个分支中，因而含于单个切片 $S_{i_j}\subseteq W_{i_j}$ 中。对每个 $j=1,\dots,m-1$，切片 $S_{i_j}$ 与 $S_{i_{j+1}}$ 有公共点 $\gamma(t_j)$，故切片 $S$ 从 $p_0$ 可达。

这表明含 $N$ 中某点的、某个 $W_k$ 的每个切片都从 $p_0$ 可达。为完成第二可数性的证明，只需注意每个 $S_{i_j}$ 本身都是一个积分流形，因而由命题\cref{proposition:ch19-局部结构}，它至多与 $W_{i_{j+1}}$ 的可数多个切片相交；从而从 $p_0$ 可达的切片只有可数多个。因此 $N$ 是一个 $k$ 维拓扑流形。它连通，因为它是一族有公共点的连通子空间的并。

为在 $N$ 上构造光滑结构，我们定义一个图册，它由所有形如 $(S\cap N,\psi)$ 的坐标图组成，其中 $S$ 是某个平坦坐标图的单个切片，$\psi\colon S\cap N\to\mathbb{R}^k$ 是在平坦坐标中到前 $k$ 个坐标的投影：$\psi(x^1,\dots,x^k,x^{k+1},\dots,x^n)=(x^1,\dots,x^k)$。由于任意切片都是嵌入子流形，其光滑结构唯一确定，从而当两个这样的切片 $S,S'$ 重叠时，转移映射 $\psi'\circ\psi^{-1}$ 是光滑的。关于这个光滑结构，包含映射 $N\hookrightarrow M$ 是一个光滑浸入（因为它在每个切片上都是光滑嵌入），且在每一点 $q\in N$ 处 $N$ 的切空间等于 $D_q$（因为对切片而言这是成立的）。
\end{proof}

\begin{figure}[ht]
\centering
\includegraphics[width=0.8\textwidth]{fig-19-7.png}
\caption{积分流形的并}\label{figure:ch19-union}
\end{figure}

\begin{figure}[ht]
\centering
\includegraphics[width=0.8\textwidth]{fig-19-8.png}
\caption{从 $p_0$ 可达的切片 $S$}\label{figure:ch19-accessible}
\end{figure}

\textbf{整体 Frobenius 定理的证明。} 对每个 $p\in M$，令 $L_p$ 为含 $p$ 的 $D$ 的所有连通积分流形的并。由上述引理，$L_p$ 是含 $p$ 的一个连通积分流形，且它显然是极大的。若任意两个这样的极大积分流形 $L_p$ 与 $L_{p'}$ 相交，则它们的并 $L_p\cup L_{p'}$ 是同时含 $p$ 与 $p'$ 的积分流形，故由极大性 $L_p=L_{p'}$。因而各个极大连通积分流形要么互不相交，要么相同。

若 $(U,\varphi)$ 是 $D$ 的任意平坦坐标图，则由命题\cref{proposition:ch19-局部结构}，$L_p\cap U$ 是切片的可数多个开子集的并。对任意这样的切片 $S$，若 $L_p\cap S$ 既非空又非整个 $S$，则 $L_p\cup S$ 是真正含 $L_p$ 的连通积分流形，这与 $L_p$ 的极大性矛盾。因此 $L_p\cap U$ 恰好是切片的可数并，故族 $\{L_p\colon p\in M\}$ 就是所求的叶状结构。

设 $M$ 是光滑流形，$\Phi\colon M\to M$ 是微分同胚。若 $\mathrm{d}\Phi(D)=D$，或更确切地说，若对每个 $x\in M$ 都有 $\mathrm{d}\Phi_x(D_x)=D_{\Phi(x)}$，则称 $M$ 上的分布 $D$ 是 $\Phi$-不变的。类似地，若对 $\mathscr F$ 的每一叶 $L$，子流形 $\Phi(L)$ 也是 $\mathscr F$ 的一叶，则称 $M$ 上的叶状结构 $\mathscr F$ 是 $\Phi$-不变的。

\begin{proposition}\label{proposition:ch19-不变分布与叶状}
设 $M$ 是光滑流形，$\Phi\colon M\to M$ 是微分同胚。设 $D$ 是 $M$ 上的一个对合分布，$\mathscr F$ 是它所确定的叶状结构。则 $D$ 是 $\Phi$-不变的当且仅当 $\mathscr F$ 是 $\Phi$-不变的。
\end{proposition}

\begin{proof}
习题 19-9。
\end{proof}
\end{document}
